\section{Fourteenth Century through Lateran V: Poverty, Authority, and Reform}

The late medieval record is unusually rich in decretals, conciliar acts, university determinations, inquisitorial files, and civil statutes. That abundance does not make its labels simple. Mystical propositions were sometimes attributed to a supposed sect on the basis of hostile reports; evangelical-poverty disputes moved between discipline and doctrine; reform movements divided internally; and ecclesiastical judgments often preceded distinct civil punishments.

\begin{heresydossier}{The Templars: allegations, suppression, and the absence of a definitive heresy judgment}{dossier:templars}
\objectfield Philip IV of France arrested the Knights Templar in 1307 and promoted accusations that admission rites required denial of Christ, spitting on the cross, illicit kisses, permission of sodomy, worship of a head or idol, and unauthorized absolution. Records contain admissions, denials, retractions, and regional acquittals; torture and coercive custody materially compromise many confessions. Evidence grade A for the proceedings, C for a uniform secret creed.

\propositionfield The alleged acts, if freely performed and believed, could involve apostasy, sacrilege, sexual wrongdoing, superstition, or usurpation of sacramental power. The evidence does not establish that the international order corporately professed one coherent heresy, and the charges should not be converted into propositions the Templars are known to have taught.

\responsefield Clement V ordered a church-wide inquiry after the French arrests. At Vienne, the bull \emph{Vox in excelso} of 22 March 1312 suppressed the order by apostolic provision because scandal and suspicion made its continuation impracticable; it expressly did not proceed by a definitive judicial sentence establishing the order's guilt. \emph{Ad providam} transferred most property to the Hospitallers, while individual cases were reserved or returned to competent tribunals.

\aftermathfield Philip's officers had initiated imprisonment and torture, and French secular authority burned fifty-four Templars in 1310 and the relapsed leaders Jacques de Molay and Geoffroi de Charney in 1314. Those executions were not a conciliar definition. The case belongs in the census because it became a paradigmatic heresy prosecution, but its historically specific conclusion is suppression amid unproved corporate allegations, not a Templar heresy with settled articles.
\end{heresydossier}

\begin{heresydossier}{Vienne on the rational soul as form of the body}{dossier:vienne-soul-form}
\objectfield Medieval theologians disputed how the intellectual soul, bodily life, and subordinate vital powers constitute one human being. The Council of Vienne's decree is anonymous; although its history is often connected with propositions attributed to Peter John Olivi, the act names no author or organized sect. Evidence grade A for the definition, B/C for its intended personal target.

\propositionfield The delimited error asserts or doubts that the substance of the rational or intellectual soul is of itself and essentially the form of the human body. The definition protects the substantial unity of a true human body and rational soul; it does not by itself decide every scholastic question about subordinate forms, faculties, individuation, or embryology.

\responsefield Vienne's first doctrinal decree, approved in 1312 and later promulgated in the \emph{Clementines}, rejected the contrary proposition as erroneous and opposed to Catholic truth. It defined that anyone who thereafter stubbornly asserts, defends, or holds it is to be considered a heretic. The personal classification therefore includes the act's express condition of future obstinate adherence.

\aftermathfield The definition became a controlling premise of Catholic anthropology and was later recalled by Lateran V. No historical evidence warrants inventing a Vienne soul-form sect or declaring every theologian whose technical account differs from Thomism a condemned heretic.
\end{heresydossier}

\begin{heresydossier}{Beguards, Beguines, and propositions attributed to the Free Spirit}{dossier:beguards-beguines}
\objectfield Beguines and Beguards formed diverse, usually local communities of lay religious life without one rule or central government. Alongside orthodox houses, inquisitors and theologians reported antinomian or quietist teachings under the label \emph{Free Spirit}. Marguerite Porete's \emph{Mirror of Simple Souls}, condemned at Paris before her execution in 1310, is an identifiable text; it is not a creed for every Beguine. Evidence grade A for the conciliar text, B/C for the extent and membership of a sect.

\propositionfield The eight errors recited in the Council of Vienne decree \emph{Ad nostrum} include claims that a person can reach such perfection as to become impeccable; that the perfected need neither fast nor pray nor obey ordinary moral commands; that sensual acts cannot be sinful when nature seeks them; that the perfected soul need not reverence the Eucharist; and that beatific vision is naturally attainable in this life. These are condemned propositions recorded by the council, not demonstrated beliefs of all people called Beguines or Beguards.

\responsefield Vienne in 1311--1312 condemned the eight propositions and those who knowingly and obstinately maintained them; John XXII promulgated the decree in the \emph{Clementines} in 1317. The companion decree \emph{Cum de quibusdam mulieribus} restricted an unapproved Beguine status while expressly leaving room for faithful women who lived honorably in penance. Local investigations still had to establish what a particular person had actually held.

\aftermathfield Some communities were suppressed, some were examined and allowed to continue, and others were regularized or affiliated with approved orders. \emph{Free Spirit} remained a mobile polemical and judicial category. Neither the execution of Porete by secular authority nor Vienne's censure licenses a collective judgment on medieval laywomen's spirituality.
\end{heresydossier}

\begin{heresydossier}{Spiritual Franciscans and the Fraticelli}{dossier:spiritual-franciscans}
\objectfield Disputes within the Franciscan family concerned observance of poverty, papal regulation of the rule, and apocalyptic claims about the order and Church. The terms \emph{Spirituals}, \emph{fraticelli}, and the many groups called \emph{bizochi} were overlapping but not interchangeable. Papal constitutions document the acts well; much evidence about dissident networks comes from prosecutors. Evidence grade A/B.

\propositionfield The first conflict often concerned refusal of commands about habit, storehouses, or community life, which was disciplinary. Some separatists then claimed that the carnal Roman Church had forfeited authority and that their community alone was the true spiritual Church. A later, precisely framed dispute concerned the assertion that Christ and the apostles possessed nothing, either individually or in common, and that this total absence of ownership was a doctrine of Scripture.

\responsefield John XXII's \emph{Quorumdam exigit} in 1317 commanded obedience on disputed observances; \emph{Sancta Romana} later that year suppressed unauthorized communities; and \emph{Gloriosam Ecclesiam} in 1318 condemned errors attributed to an obstinate Sicilian group. In \emph{Cum inter nonnullos} in 1323, John declared the categorical no-ownership assertion erroneous and heretical. The acts therefore range from governance of an order to censure of specified ecclesiological and poverty propositions.

\aftermathfield Four Spiritual Franciscans were burned by civil authority at Marseille in 1318 after refusing submission, and later Fraticelli communities faced inquisitorial and secular suppression. Other strict Franciscans remained or were reconciled within the Church, eventually contributing to observant reform. Zeal for poverty as such was not condemned; the decisive issues were the proposition asserted and, in some cases, contumacious separation.
\end{heresydossier}

\begin{heresydossier}{Marsilius of Padua, John of Jandun, and ecclesial coercive power}{dossier:marsilius-jandun}
\objectfield Marsilius's \emph{Defensor pacis}, written in 1324 with assistance or association traditionally assigned to John of Jandun, joined a theory of political sovereignty to claims about the Church. Both men entered the protection of Louis of Bavaria during his conflict with John XXII. The book and papal act survive. Evidence grade A.

\propositionfield The papal summary selected five conclusions: Christ's payment of tribute was treated as compelled subjection to temporal power; Peter was denied any headship over the other apostles and Christ any appointment of a vicar; the emperor was said to possess authority to appoint, punish, or depose a pope; all priests were made equal in jurisdiction by divine institution; and coercive punishment of religious error was reserved to the civil legislator. The larger work also located supreme legislative authority in the body of citizens and treated ecclesiastical law as lacking coercive force without civil enactment.

\responsefield John XXII's constitution \emph{Licet iuxta doctrinam}, 23 October 1327, extracted and condemned the five propositions as contrary to Scripture and Catholic faith, with differentiated theological reasoning. The censure arose within a papal-imperial struggle, but it addressed identifiable claims about apostolic primacy, orders, jurisdiction, and the Church's authority rather than merely Marsilius's political allegiance.

\aftermathfield Marsilius was never brought to a completed personal trial at Avignon and died under imperial protection. His political ecclesiology influenced later conciliar and state-sovereignty arguments, but that intellectual reception does not make every advocate of representative government a Marsilian. Imperial campaigns and penalties for treason or rebellion belong to the political history, not to the content of the doctrinal censure.
\end{heresydossier}

\begin{heresydossier}{Meister Eckhart's twenty-eight propositions}{dossier:eckhart-propositions}
\objectfield Dominican theologian Meister Eckhart faced an inquisitorial inquiry at Cologne and appealed to Avignon. His Latin and German works, defense records, the Avignon report, and John XXII's bull preserve an unusually rich but textually difficult dossier. Extraction from sermons can remove paradox, qualification, or rhetorical setting. Evidence grade A/B.

\propositionfield The twenty-eight articles concern the relation of creation to God's eternity, claims that even evil acts manifest divine glory, detachment from created goods and reward, the deified person's relation to the Son and divine action, and an uncreated or uncreatable element in the soul. The bull itself acknowledges that eleven articles could bear a Catholic sense with explanations; no single pantheist slogan accurately represents all twenty-eight.

\responsefield John XXII's \emph{In agro dominico}, 27 March 1329, condemned the first fifteen and last two articles as heretical and the intervening eleven as evil-sounding, rash, and suspect of heresy. It directed proceedings against defenders according to those distinct grades rather than distributing one censure across all. The bull also recorded Eckhart's prior Catholic profession and revocation of anything in his preaching or writings capable of an heretical or erroneous sense.

\aftermathfield Eckhart died before the bull and was not personally declared an obstinate heretic by it. His writings influenced later Catholic mysticism and are studied with renewed attention to context, while the proposition-level censures remain part of the record. Later reception cannot either erase the act or extend it to every expression of mystical union and deification.
\end{heresydossier}

\begin{heresydossier}{The Black Death flagellants}{dossier:flagellants}
\objectfield Organized flagellant companies spread through central and northern Europe during the plague of 1348--1350. Public bodily penance had earlier orthodox precedents, but some companies followed unauthorized rituals, circulated a purported heavenly letter, claimed powers normally connected with sacramental ministry, and acted independently of bishops. Chronicles and the papal prohibition are principal witnesses. Evidence grade A/B, with considerable local variation.

\propositionfield Reports attribute salvific necessity or quasi-sacramental efficacy to a fixed course of scourging, lay absolution, contempt for clergy and sacraments, and the claim that the movement's letter conveyed a new divine mandate. No evidence warrants attributing every claim to every penitent. Voluntary self-discipline or public prayer, considered by itself, was not the condemned object.

\responsefield Clement VI's \emph{Inter sollicitudines}, 20 October 1349, ordered the unauthorized associations dissolved and directed bishops and secular rulers to restrain them, while permitting properly governed penitential practice. The response defended episcopal oversight, sacramental confession, and the sufficiency of public revelation against a self-authorizing movement.

\aftermathfield Suppression varied by territory and included civil arrest and execution. Flagellant crowds were also entangled in anti-Jewish violence during the plague, but responsibility must be assigned to documented actors rather than to every procession. Later Catholic confraternities of discipline operated under ecclesiastical rule and are not continuations of the censured claims.
\end{heresydossier}

\begin{heresydossier}{John Wyclif}{dossier:wyclif}
\objectfield The Oxford theologian John Wyclif developed positions on dominion, Eucharist, predestination, Scripture, and ecclesiastical office during the papal schism and English crown-clergy conflict. His own Latin and English works survive alongside proposition lists whose wording sometimes compresses or alters their context. Evidence grade A/B.

\propositionfield Censured propositions include the persistence of bread and wine after consecration and rejection of transubstantiation; an account of the Church primarily as the congregation of the predestined; claims that mortal sin or lack of grace defeats a prelate's true dominion; denial or radical restriction of papal jurisdiction; and assertions permitting temporal lords to remove church goods. Not every list gave each item the same grade, and Wyclif's appeals to Scripture cannot be reduced to a simple denial of all tradition.

\responsefield Gregory XI's five letters of 22 May 1377 transmitted nineteen propositions for examination. The Blackfriars council at London in 1382 judged ten conclusions heretical and fourteen erroneous. The Council of Constance, session 8 on 4 May 1415, condemned forty-five articles and Wyclif's books, while distinguishing articles already condemned as heretical from others censured as erroneous, scandalous, or seditious; Martin V's \emph{Inter cunctas} in 1418 supplied interrogatories for adherents. These are several acts with mixed censures, not one undifferentiated anathema.

\aftermathfield Wyclif died in communion neither tried nor executed as a living convicted heretic. Constance ordered posthumous removal of his remains, carried out in 1428. His works influenced English Lollards and Bohemian reformers, but both receptions selected and altered his positions. English political restrictions and later civil penalties require separate attribution.
\end{heresydossier}

\begin{heresydossier}{Lollardy as a diverse English reception}{dossier:lollardy}
\objectfield The label \emph{Lollard} named networks of clerics, gentry, artisans, and readers influenced to different degrees by Wyclif from the 1380s onward. Trial records, manuscript books, episcopal registers, and the anonymous \emph{Twelve Conclusions} presented to Parliament in 1395 show overlap without a single membership roll or confession. Evidence grade A/B.

\propositionfield Recurrent positions questioned transubstantiation, clerical celibacy, confession to priests, images and pilgrimage, prayers for the dead, papal authority, and the legitimacy of war or capital punishment. The \emph{Twelve Conclusions} joined several of these to a program of church disendowment, but it cannot be assigned to every accused Lollard. Vernacular Scripture itself was not a heresy; unauthorized translation, teaching, and the doctrinal use made of texts were separately contested.

\responsefield English bishops and the universities examined particular suspects and propositions under existing canon law. Archbishop Thomas Arundel's provincial constitutions, issued in 1409 after the Oxford deliberations, regulated vernacular biblical translation, preaching, schools, and theological teaching. They were a local disciplinary and preventive regime, not an ecumenical list of Lollard dogmas.

\aftermathfield Parliament and the crown added the statute \emph{De heretico comburendo} in 1401, authorizing secular burning after ecclesiastical conviction; later confiscations and executions were therefore mixed church-state proceedings. Recantation, penance, underground survival, and renewed activity during the Reformation all occurred. Accusation records must not be read as proof that a lay reader held Wyclif's complete theology.
\end{heresydossier}

\begin{heresydossier}{Jan Hus and the Council of Constance}{dossier:hus}
\objectfield Jan Hus, a priest and preacher at Prague, combined Czech reform concerns with selective reception of Wyclif. His sermons, treatise \emph{On the Church}, letters, trial record, and Constance's thirty articles permit comparison between his words and the council's judicial extracts. He denied having taught some formulations and disputed the condemned meaning of others. Evidence grade A/B.

\propositionfield The articles concerned the Church as the body of the predestined; whether a sinful or reprobate pope remains a true member and worthy minister of the Church; the conditions of ecclesiastical obedience and excommunication; the authority of morally corrupt prelates; and Hus's defense of censured Wycliffite propositions. Some statements could bear an orthodox qualified sense, while others were judged to undermine visible office and jurisdiction; attribution and sense were part of the trial dispute.

\responsefield Constance's session 15 on 6 July 1415 condemned thirty articles, censured Hus's appeal and persistence, and declared him obstinate after he refused the demanded abjuration. Hus maintained that he could not recant errors he had never held and would retract anything shown false from Scripture; the tribunal did not accept that limitation. The council degraded him from clerical status and relinquished him to secular authority.

\aftermathfield Secular officials of King Sigismund burned Hus the same day; the council did not itself carry out the execution. His death intensified a Bohemian religious and national crisis and made him a martyr to Hussite traditions. A historical account can state both the council's judgment and Hus's defense without retroactively declaring every Czech reform demand heretical.
\end{heresydossier}

\begin{heresydossier}{Hussite and Utraquist currents}{dossier:hussites-utraquists}
\objectfield After Hus's death, Bohemian reform coalesced around the Four Articles of Prague: freer preaching, communion from the chalice for the laity, correction of public mortal sin, and limitation of clerical temporal power. Utraquists, Taborites, Orphans, and later Unitas Fratrum communities differed sharply in sacramental theology, ministry, and politics. Evidence grade A/B.

\propositionfield The central Utraquist claim made reception under both kinds a divine command or necessary practice for the laity. Radical Taborite teachers variously denied transubstantiation or the real presence, rejected sacramental forms and purgatory, and announced apocalyptic transformation; moderates did not share all such propositions. The chalice itself was not intrinsically heretical, since Christ is whole under either species and the Church has at times permitted both.

\responsefield Constance session 13, 15 June 1415, taught the sufficiency of communion under either species and condemned as erroneous the unauthorized claim and practice that the chalice was necessary for the laity. Papal and imperial crusades against Bohemia were military measures distinct from that doctrinal decree. Negotiations with the Council of Basel produced the Compactata of 1436, permitting the chalice under specified Catholic teaching and discipline; Pius II later declared the settlement no longer binding in 1462.

\aftermathfield The Bohemian wars ended neither by simple extermination nor wholesale Roman acceptance: Utraquist worship received a limited legal settlement, radical armies were defeated at Lipany in 1434, and communities developed along several confessional lines. Later permissions for communion under both kinds confirm that mode of reception and the doctrine attached to it must be distinguished.
\end{heresydossier}

\begin{heresydossier}{Conciliarism and appeals from pope to council}{dossier:conciliarism}
\objectfield The Great Western Schism made a general council the practical means of restoring one recognized pope and generated competing theories of ecclesial representation. Conciliarism in the strict sense held that a general council possesses authority over a pope, at least in matters of faith, schism, and reform; advocates differed over whether that superiority was exceptional or permanent. Evidence grade A, with disputed juridical reception of individual acts.

\propositionfield The contested propositions were not the legitimacy of councils, episcopal collegiality, or reform through synodal deliberation. They concerned whether a council derives jurisdiction immediately from Christ in such a way that every Christian including a reigning pope is juridically bound to obey it, and whether an individual may evade a papal judgment by appealing to a future general council.

\responsefield Constance session 5 issued \emph{Haec sancta} in 1415 during the schism, and session 39 issued \emph{Frequens} in 1417 to require regular councils. The authority, scope, and papal confirmation of \emph{Haec sancta} remain disputed in Catholic interpretation. During the later Basel conflict, Eugene IV rejected the assembly's claim to judge and depose him. Pius II's \emph{Execrabilis} in 1460 definitively prohibited appeals from a pope to a future council as an abuse contrary to ordered jurisdiction.

\aftermathfield Constance successfully ended the schism, and councils remain institutions of the Church; the censure concerns a supremacy claim, not conciliar deliberation itself. Basel's rump elected antipope Felix V, who later resigned, while reunion work continued at Ferrara--Florence under papal presidency. Lateran V's \emph{Pastor aeternus} in 1516 supplied a later authoritative rejection of conciliar superiority, but it belongs chronologically to the next century's reception.
\end{heresydossier}

\begin{heresydossier}{Lateran V on the individual soul's immortality}{dossier:lateran-v-soul}
\objectfield Late-medieval and Renaissance Aristotelian debates asked whether philosophy could establish personal immortality and whether one separate intellect served all human beings. Lateran V addressed propositions, not a named Averroist church, and its decree predates Pietro Pomponazzi's best-known intervention. Evidence grade A.

\propositionfield The targeted alternatives say that the rational soul is mortal by nature or that one and the same soul or intellect animates all human bodies. The council also rejected a two-truth defense in which a proposition could be true in philosophy while contradicting revealed truth; it did not condemn philosophical inquiry or every dispute about the demonstrability of immortality.

\responsefield Leo X promulgated \emph{Apostolici regiminis} with the approval of Lateran V in session 8, 19 December 1513. Recalling Vienne, it taught that the rational soul is genuinely, of itself, and essentially the body's form, immortal, and multiplied according to the multitude of bodies, and condemned contrary assertions as false and heretical. It also required philosophy teachers to answer such arguments and clergy studying philosophy to complete theological formation.

\aftermathfield Debate continued over what natural reason could prove, especially after Pomponazzi's 1516 treatise, even among thinkers who professed the doctrine by faith. The decree settled the Catholic propositions of personal immortality and individual souls; it did not name every Aristotelian interpreter a formal heretic.
\end{heresydossier}
